\chapter{Macro System}
\label{ch:macros}

% Source: src/zmacro.c

\section{Overview}

Zinc macros are syntactic pattern-matching rules that expand to arbitrary AST fragments at compile time.
They are declared with the \zn{macro} keyword and expanded during the parsing phase (before type checking).

\begin{lstlisting}[style=zinc, caption={A simple macro declaration}]
  macro unless $cond $body -> {
      if !($cond) {
          $body
      }
  }
\end{lstlisting}

\begin{lstlisting}[style=zinc, caption={Using the macro}]
  unless is_valid {
      return
  }
\end{lstlisting}

\section{Pattern Elements}

A macro pattern is a sequence of \emph{pattern elements}:

\begin{center}
\begin{tabular}{lp{9cm}}
\toprule
\textbf{Syntax} & \textbf{Meaning} \\
\midrule
literal token        & Must match exactly (e.g.\ \zn{from}, \zn{to}). \\
\texttt{\$name}      & Capture variable: matches a single expression or statement block. \\
\texttt{@name}       & Identifier capture: matches one identifier and introduces it as a new name. \\
\texttt{\#name}      & Type capture: matches a type expression. \\
\texttt{'\,tok}      & Quoted token: forces a literal match on an otherwise overridable symbol. \\
\texttt{\$(...)*}      & Kleene-star repetition of the enclosed sub-pattern. \\
\texttt{\$(...)+}      & One-or-more repetition. \\
\bottomrule
\end{tabular}
\end{center}

\section{Expansion}

During parsing, when the parser encounters a token sequence that begins with a macro keyword, \texttt{zmacro.c} is invoked to:

\begin{enumerate}
  \item Match the remaining tokens against the declared pattern.
  \item Bind each capture variable to the matched token span.
  \item Substitute the captures into the macro body, producing a new token stream.
  \item Push the new token stream back into the parser's \texttt{ZTokenStream} stack so it is parsed as ordinary Zinc code.
\end{enumerate}

The \texttt{ZTokenStream} maintains a \texttt{prev} pointer, forming a stack of streams.
Once the injected stream is exhausted, parsing continues from the outer stream.

\section{Hygiene}

Identifier captures (\texttt{@name}) introduce fresh names that do not collide with identifiers already in scope.

% TODO: describe the gensym mechanism in zmacro.c

\section{Example: Custom Loop}

\begin{lstlisting}[style=zinc, caption={A range-based \texttt{for} macro}]
  macro for @item from $start to $end $body -> {
      @item := $start
      for @item < $end {
          $body
          @item = @item + 1
      }
  }
\end{lstlisting}

\begin{lstlisting}[style=zinc, caption={Usage}]
  for i from 0 to 10 {
      printf("%d\n", i)
  }
\end{lstlisting}

\section{Macros and \texttt{goto}}

Because expansion produces ordinary Zinc source, macros can use \zn{goto} and labels internally to express complex control structures, such as break-from-nested-loop or custom early-exit patterns.

\begin{lstlisting}[style=zinc, caption={A \texttt{while} macro using \texttt{goto}}]
  macro while $cond $body -> {
      __entry:
      if !($cond) { goto __exit }
      $body
      goto __entry
      __exit:
  }
\end{lstlisting}

\begin{warn}
  Because macro labels are emitted literally, two invocations of the same goto-using macro in the same function will produce duplicate label definitions.
  The \texttt{@name} identifier capture should be used to introduce fresh label names.
\end{warn}
